\chapter{Anexos}

\section{Datasheets y Documentación Técnica}

\subsection{Motor D5065 270KV}
\begin{itemize}
    \item KV Rating: 270 RPM/V
    \item Pole Pairs: 7
    \item Torque Constant: 8.27/270 Nm/A ($\approx$0.0306 Nm/A)
    \item Max Current: 30A (configurado)
    \item Encoder: Integrado 8192 CPR
\end{itemize}

\subsection{Encoder AS5600}
\begin{itemize}
    \item Tipo: Magnetic rotary encoder
    \item Resolución: 12-bit (4096 positions/rev)
    \item Interfaz: I2C (dirección 0x36)
    \item Voltaje: 3.3V - 5V
\end{itemize}
Datasheet: \url{https://ams.com/as5600}

\subsection{Multiplexor I2C TCA9548A}
\begin{itemize}
    \item Canales: 8 buses I2C independientes
    \item Dirección I2C: 0x70 (configurable)
    \item Voltaje: 1.65V - 5.5V
\end{itemize}
Datasheet: \url{https://www.ti.com/product/TCA9548A}

\subsection{ODrive v3.6}
\begin{itemize}
    \item Firmware: v0.5.6 (oficial)
    \item Comunicación: USB, UART
    \item Motores soportados: 2 (BLDC/PMSM)
    \item Corriente máxima: 48A (24V), 24A (56V)
\end{itemize}
Documentación: \url{https://docs.odriverobotics.com/}

\section{Código Fuente}

El código fuente completo del proyecto está disponible en:
\begin{center}
\url{https://github.com/Josama-99/tfg_biped_leg}
\end{center}

\subsection{Estructura del Repositorio}
\begin{verbatim}
tfg_biped_leg/
├── config/                     # Configuraciones de motores
│   └── my_config.json
├── launch/                     # Archivos launch de ROS2
├── msg/                        # Mensajes ROS2
├── srv/                        # Servicios ROS2
├── src/                        # Drivers de bajo nivel
│   ├── as5600_encoder.py       # Driver AS5600
│   └── i2c_mux.py             # Driver TCA9548A
├── tfg_biped_leg/             # Paquete principal Python
│   ├── odrive_interface.py     # Control ODrive
│   ├── odrive_enums.py         # Enums ODrive
│   ├── odrive_driver.py        # Driver ROS2
│   ├── leg_kinematics.py       # Cinemática
│   └── trajectory_generator.py  # Trayectorias
├── scripts/                    # Scripts de utilidad
│   ├── test_single_motor.py
│   ├── calibrate_motors.py
│   └── odrive_test.py
└── tests/                      # Tests unitarios
\end{verbatim}

\subsection{Ejemplo de Código: Conexión ODrive}
\begin{lstlisting}[language=Python, caption=Conexión y configuración básica de ODrive]
import odrive
from odrive.enums import *

def connect_odrive():
    """Conectar y configurar ODrive"""
    print("Buscando ODrive...")
    odrv = odrive.find_any()
    print(f"ODrive encontrado: S/N {odrv.serial_number}")
    return odrv

def configure_motor(odrv, axis_num):
    """Configurar motor en eje especificado"""
    axis = getattr(odrv, f'axis{axis_num}')
    axis.motor.config.current_lim = 30
    axis.motor.config.calibration_current = 10
    axis.encoder.config.cpr = 8192
    print(f"Motor en eje {axis_num} configurado")

# Uso
odrv = connect_odrive()
configure_motor(odrv, 0)
\end{lstlisting}

\section{Diagramas}

\subsection{Diagrama de Arquitectura del Sistema}
El diagrama de arquitectura (Figura \ref{fig:arquitectura} en el Capítulo 3) muestra la interconexión entre:
\begin{itemize}
    \item Raspberry Pi (control de alto nivel)
    \item ODrive Oficial y Makerbase (control de motores)
    \item Multiplexor TCA9548A (lectura de encoders)
    \item Motores y encoders
\end{itemize}

\subsection{Esquemático de Conexiones}
\begin{itemize}
    \item \textbf{ODrive - Motor}: Fases U, V, W; Encoder A, B
    \item \textbf{Raspberry Pi - ODrive}: USB (comunicación serie)
    \item \textbf{Raspberry Pi - TCA9548A}: I2C (SDA, SCL)
    \item \textbf{TCA9548A - AS5600}: Canales 0, 1, 2 (un encoder por canal)
\end{itemize}

\section{Configuración del Entorno de Desarrollo}

\subsection{Raspberry Pi Setup}
\begin{enumerate}
    \item Instalar ROS2 Jazzy en Ubuntu 24.04
    \item Crear entorno virtual Python: \texttt{python3 -m venv tfg\_venv}
    \item Instalar dependencias: \texttt{pip install odrive numpy smbus2 pyyaml colorama matplotlib}
    \item Configurar permisos USB para ODrive (regla udev)
    \item Habilitar I2C: \texttt{sudo raspi-config → Interface Options → I2C}
\end{enumerate}

\subsection{Repositorio de Referencia}
Para el código de control de motores, se utilizó como referencia el repositorio privado:
\begin{center}
\url{https://github.com/SergiMuac/odrive_smapy}
\end{center}
(Disponible bajo solicitud)
